\setcounter{section}{0}
\setcounter{subsection}{0}
\setcounter{subsubsection}{0}

\clearpage
\begin{center}
  {\Large\bfseries In-Class Lab}\\[0.5em]
\end{center}
\vspace{1em}

\section{In-Class Lab: $3$-Version Development and Independence Assessment}

In the in-class portion of this lab, you will construct three independently generated implementations of the tax engine specified in~\cref{appendix:tax-spec}.
The objective is to instantiate the $2$-out-of-$3$ ($2$oo$3$) voting architecture analyzed in the pre-lab.
By the end of the session, each team will have produced three distinct implementations of the same frozen specification and evaluated them on a visible test set.

To approximate independent development:

\begin{itemize}
  \item Each implementation must be generated in a fresh LLM session.
  \item Team members may discuss high-level strategy before generation.
  \item During generation, team members may not share code across versions.
  \item Each implementation should intentionally vary along at least one dimension (\eg language, architecture, numeric handling, library usage).
\end{itemize}

\subsection{Team Structure and Requirement}

All submissions (in-class and online) must contain three LLM-generated implementations.

\textbf{In-class students:}

Work in groups of two students.
At most one group of three may be formed if necessary.

\textbf{Online students:}

Online students should work with those of their Team project partners who are also online students.
This will result in groups of 2-3 online students for this submission.

\begin{tcolorbox}[colback=purple!10,colframe=purple!70!black,boxrule=0.5pt]
\textbf{Team members:}  
\begin{itemize}
    \item Ming Jia (jia213@purdue.edu)
    \item Dovid Korn (korn4@purdue.edu)
\end{itemize}
\textbf{Team signup forum:} \labTwoForumLink\par
Please indicate your team members in the Brightspace discussion forum link above.
\end{tcolorbox}

\subsection{Version Independence Requirements}



You may use prompting strategies such as those illustrated in \cref{appendix:sample-prompts}.
However, the specification in \cref{appendix:tax-spec} must remain unchanged.

\subsection{Plan for Diversity}

Before generating code, discuss how you will induce diversity across the three implementations.
You may draw on prompting strategies such as those illustrated in~\cref{appendix:sample-prompts}.


\begin{itemize}
  \item How you will induce diversity across the three implementations,
  \item What dimensions you expect might influence failure behavior,
  \item Whether floating-point handling, architecture, or library choice might affect correctness.
\end{itemize}

\begin{tcolorbox}[colback=purple!10,colframe=purple!70!black,boxrule=0.5pt]
\textbf{Team Strategy:}  
We induce diversity by implementing the same specification in three different programming languages: Python (impl-1), Go (impl-2), and C++ (impl-3).
Using distinct languages yields different runtime behavior, standard libraries, and numeric handling, which may influence failure modes and the independence of faults across versions.
\end{tcolorbox}

\subsection{Generate Three Implementations}

Generate your three implementations.

\begin{enumerate}
  \item Start a new LLM session.
  \item Provide the full specification from \cref{appendix:tax-spec}.
  \item Develop a complete implementation that reads the input and produces tax outputs.
  \item Ensure the implementation runs locally.
\end{enumerate}

\textit{\textbf{Note 1:}
I have fed the specification to ChatGPT, Microsoft CoPilot, and Google Gemini and prompted them to ``Write a program to answer this spec''.
All of them choked, had network issues, or timed out.
For the in-class part, you will need to break things down for them so you can complete the implementation in a timely manner.}

\textit{\textbf{Note 2:}
I advise you to complete two implementations and do a 1oo2 check before you develop the third implementation.
}

A small smoke test set (Test Set 1) is provided on Brightspace.
Because we do not provide a reference implementation for this lab, Test Set 1 cannot help you establish correctness (nor is that the point of this lab).
Instead, it is used to reduce noise by ensuring that all implementations:

\begin{itemize}
  \item parse inputs successfully,
  \item produce well-formed outputs, and
  \item behave consistently enough to proceed with meaningful analysis of disagreement and potential correlated failures.

\end{itemize}

Run all three implementations on Test Set 1.
If any implementation crashes, produces malformed output, or disagrees with the other implementations, revise it (\eg via debugging and prompting) and re-run Test Set 1.
Proceed only once your team has a stable outcome on Test Set 1.
You should aim for unanimous agreement across all cases, thus indicating each version is (seemingly) correct.
Settling for 2oo3 agreement on every case would suffice if needed.

\subsection{Independence Analysis on Test Set 2}

A second, larger ``full'' test set (Test Set 2, approximately 10{,}000 inputs) is provided on Brightspace.
A ``runner'' program is also supplied to let you drive multiple implementations and get results.
Run all three implementations on Test Set 2 and record, for each input, whether the three implementations agree.

Because no external oracle is provided, we treat disagreement among implementations as evidence of potential fault events.
This allows us to study the structure of divergence without assuming any version is correct.

Construct an agreement table with one row per input from Test Set 2.
For each input, record the outputs produced by each version:

\[
\text{Input} \rightarrow (\text{Output}_{V1},\ \text{Output}_{V2},\ \text{Output}_{V3})
\]


Using this data, compute and report:

\begin{itemize}
  \item The fraction of inputs with unanimous agreement (3/3).
  \item The fraction of inputs with exactly one dissenting version (a 2--1 split).
  \item The fraction of inputs where multiple versions diverge from one another.
  \item For each version, how often it disagrees with the other two.
\end{itemize}

Discuss:

\begin{itemize}
  \item Whether disagreement events appear isolated or clustered on specific inputs.
  \item Whether the same implementation tends to disagree repeatedly.
  \item Whether disagreement patterns suggest independence or correlated behavior.
\end{itemize}

At this stage, you are not estimating true failure probabilities.
Rather, you are examining the empirical structure of disagreement as preliminary evidence about independence assumptions.

This analysis evaluates \emph{consistency} and \emph{correlated divergence} among independently generated implementations.
It does not, by itself, establish correctness, since we have not included correct answers in the test set.

\subsection{Preliminary Analysis}

Within your team, discuss:

\begin{itemize}
  \item Are disagreement events concentrated on the same inputs across implementations, or do they appear scattered?
  \item When disagreements occur, do they tend to implicate the same implementation repeatedly (one version often being the outlier), or do outliers rotate across versions?
  \item What aspects of the specification appear to trigger disagreements?
  For example, are disagreements concentrated on inputs that activate the floating-point rules (e.g., the EWMA-based surcharge), threshold boundaries, or phaseouts?
  \item Based on the observed disagreement patterns, is the independence assumption plausible, or do you see evidence of correlated behavior (common-mode mistakes)?
\end{itemize}

You are not required to perform formal statistical testing at this stage.
The goal is to build intuition about correlated behavior and the independence assumption using empirical disagreement data.
If you choose to complete the take-home portion of the lab, you will perform a larger version of this experiment and may incorporate more formal modeling of the underlying hypothesis of independence.

\subsection{Agreement Results on Test Set 2 (All Three Implementations)}

All three implementations (impl-1 Python, impl-2 Go, impl-3 C++) were run on Test Set 2.
The agreement table CSV has columns \texttt{taxpayer\_id}, \texttt{out\_v1}, \texttt{out\_v2}, \texttt{out\_v3} (federal + state tax totals per version).

\begin{table}[h]
\centering
\begin{tabular}{lc}
\toprule
Metric & Value \\
\midrule
Total inputs (Test Set 2) & 10{,}000 \\
Unanimous agreement (3/3) & 9{,}653 (96.53\%) \\
Exactly one dissenting version (2--1 split) & 330 (3.30\%) \\
All three versions differ & 17 (0.17\%) \\
Times impl-1 disagrees with the other two & 347 \\
Times impl-2 disagrees with the other two & 347 \\
Times impl-3 disagrees with the other two & 347 \\
\bottomrule
\end{tabular}
\caption{Agreement statistics for impl-1 (Python), impl-2 (Go), and impl-3 (C++) on Test Set 2.}
\label{tab:agreement}
\end{table}

\textbf{Discussion.}
About 96.5\% of inputs show unanimous agreement; 3.3\% have a single dissenting version (2--1 split) and 0.17\% have all three outputs different.
Each implementation disagrees with the other two on 347 inputs, so no single version is consistently the outlier---disagreement is spread across implementations, which is consistent with independent (or weakly correlated) fault behavior rather than one consistently faulty version.

\textbf{Why results differ.}
Likely causes include: (1)~\textbf{Floating-point behavior}---Python, Go, and C++ use different runtime and library behavior for rounding (e.g., half-up vs.\ half-even) and for the order of operations in expressions, which can change the least significant digits in the EWMA recursion and in bracket/surcharge math; (2)~\textbf{When rounding is applied}---the spec requires rounding to the nearest dollar but does not fix intermediate rounding points, so one implementation may round after each step while another rounds only at the end, producing different totals; (3)~\textbf{Threshold and boundary cases}---comparisons such as taxable income at bracket boundaries (e.g., near \$100{,}000 or \$300{,}000) or EWMA at the surcharge threshold (\$1{,}000{,}000) are sensitive to tiny numeric differences, so one version may place a value just above the threshold and another just below; (4)~\textbf{FIFO and decimal handling}---capital gains are specified to three decimal places and then rounded; small differences in how each language rounds or accumulates over many lots can propagate into the final tax.

\begin{tcolorbox}[colback=purple!10,colframe=purple!70!black,boxrule=0.5pt]
\textbf{Deliverable for the In-Class Portion:}
Submit:
\begin{itemize}
  \item The three implementations (distinct zipped folders labeled \texttt{impl-1.zip}, \texttt{impl-2.zip}, \texttt{impl-3.zip}),
  \item Your disagreement table for Test Set 2 (as a CSV with columns \texttt{input\_id}, \texttt{out\_v1}, \texttt{out\_v2}, \texttt{out\_v3}),\footnote{For simplicity, you can sum the federal and state taxes for those outputs and just report results in the aggregate}
  \item A brief analysis describing observed agreement or disagreement patterns, including at least one table or figure.
\end{itemize}
\end{tcolorbox}